\chapter{绪论}


\section{研究背景与意义}

\subsection{研究背景}
近年来，以大语言模型（Large Language Model, LLM）为核心的生成式人工智能技术迅速发展，深刻改变了软件工程的实践方式。Hou 等人对 395 篇 LLM 软件工程（LLM4SE）文献的系统综述表明，大语言模型的应用已由早期的代码生成、代码补全与代码解释等局部任务，逐步拓展至需求工程、软件设计、软件开发、质量保障与软件维护等多个软件工程活动[1]。2025 至 2026 年的实证研究开始关注 AI 编程智能体（coding agent）在真实工程环境中的应用情况：基于 GitHub 生态数据的大规模项目分析尝试刻画 coding agent 在开源项目中的使用趋势，结果显示其已在大量真实开源项目中出现使用迹象[2]；基于 AIDev 数据集的研究则记录了 AI 辅助软件开发活动在真实仓库中的表现[3]。

然而，大语言模型在软件工程中的工程化应用仍面临突出挑战，本文将其归纳为过程可控性与上下文效率两类核心问题。

其一是过程可控性问题。大语言模型的输出具有概率性与生成式特征，以自然语言交互为主要形式的开发方式在需求澄清、任务分解、过程状态维护与结果验证等环节存在较强的过程依赖性。Li 等人的研究发现，智能体产出拉取请求的接受率与人类开发者仍存在明显差距，表明智能体的自治程度与产出可信度并非同一维度[3]--速度与自治性的提升并不天然带来工程结果的可靠性。这提示智能体的可靠性不仅取决于模型自身能力，也取决于外围工程机制，包括任务分解、状态管理、工具调用与验证反馈等环节的显式约束。因此，如何通过工程化的流程约束提高 AI 辅助开发过程的可控性与可审计性，成为值得研究的问题。

其二是上下文效率问题。随着软件仓库规模扩大，直接将大量源码、文档与开发历史注入模型上下文，会带来上下文长度受限、检索效率下降与关键信息定位困难等问题。已有研究表明，即使模型支持长上下文输入，相关信息在上下文中所处的位置仍可能显著影响模型对其的利用效果，即"Lost in the Middle"现象[4]；检索增强生成（Retrieval-Augmented Generation, RAG）研究则表明，引入外部非参数化知识库并按需检索，能够缓解参数化知识在知识访问与更新上的局限[5]。2025 至 2026 年的智能体研究进一步将仓库级上下文（repository-level context）与可执行反馈视为软件工程智能体的重要资源[6]。如何为智能体系统组织、选择与供给上下文逐渐受到关注，此类方法被统称为"上下文工程"（Context Engineering）；本文所称上下文工程，指针对智能体任务目标，对外部知识、代码结构、执行状态等信息进行组织、选择、压缩与供给的方法。

围绕上述两类核心问题，本文凝练出以下三个研究问题（Research Questions, RQ），全文各章的工作与实验均围绕这三个研究问题展开：

\begin{itemize}
\item \textbf{RQ1}：如何通过工程化的流程约束提高 AI 辅助开发过程的确定性与可控性？（对应第 4 章与创新点 1）
\item \textbf{RQ2}：如何组织项目级业务知识与代码结构知识，为智能体任务提供有效的结构化上下文？（对应第 3 章与创新点 2）
\item \textbf{RQ3}：如何在有限 Token 预算下生成任务相关、结构连贯的高密度智能体上下文？（对应第 3.3 节与创新点 3）
\end{itemize}
与上述问题相呼应，工程社区也呈现出明确的技术演进趋势：其一，规格驱动开发（Spec-Driven Development, SDD）方法受到重新关注，以 GitHub Spec Kit、OpenSpec 为代表的工具将"先明确意图与规格，再实施"逐渐确立为 AI 辅助开发的重要实践[7-8]；其二，面向智能体的工程化实践逐渐强调通过工作流、工具接口、状态管理与验证机制约束模型行为，呈现"模型能力与工程控制层协同"的设计思路；其三，以 Superpowers 为代表的技能化开源框架涌现，将工程纪律沉淀为可组合的智能体技能[9]；其四，代码知识图谱与 RAG 上下文工程兴起，为代码结构知识的组织与检索提供了新途径。然而，将上述方向统一整合、使其协同服务于 AI 辅助开发全流程的研究仍相对有限。

本文正是在上述背景下，以 Web 项目开发为典型应用场景，研究面向 AI 编程智能体的软件工程上下文构建方法--涵盖项目级知识的组织与沉淀、开发流程的工程化约束，以及任务上下文的生成与预算控制--并设计实现了 wpw（Web Project Workflow）系统作为上述方法的系统实现与实验验证平台。


\subsection{研究意义}
本文的研究意义体现在理论与实践两个层面。

在理论层面，本文为面向 AI 编程的软件工程工作流设计以及代码知识驱动的上下文生成方法提供了一种方法性探索：提出三层分离架构模式与能力-代码双层知识图谱模型，研究多源证据融合的业务-代码关联建立方法与面向大语言模型的上下文压缩方法，为相关研究方向提供可参考的工程实践案例与实验依据。

在实践层面，本文为软件工程智能体的上下文组织与生成提供了一种工程化实现路径：通过确定性流程编排降低 AI 辅助开发的使用门槛，通过知识图谱上下文提升 Token 利用效率，通过状态留痕与依赖检查提升开发过程的规范性与可追溯性，并通过能力规范的持续沉淀支持项目知识的长期积累；上述路径以 wpw 系统为载体在 Web 项目开发场景中得到端到端实现。

本课题来源于实际工程项目中 AI 辅助开发遇到的上述痛点，具有明确的问题驱动特征。


\section{国内外研究现状}
本节从 AI 辅助开发与智能体工作流编排、代码知识图谱与上下文工程两个方向综述相关研究。需要说明的是，本文涉及对比的工程工具（Spec Kit、OpenSpec、Superpowers）均为高速迭代的开源项目，其版本与生态数据检索截至 2026 年 8 月；本文对这些工具的能力描述均绑定研究采用版本，其版本号、发布日期与代码提交记录固化于附录实验环境表中以保证可复现性，不将其表述为固定的"产品属性"。


\subsection{AI 辅助开发与 Agent 工作流编排研究现状}
\textbf{（1）LLM 辅助软件开发与 Coding Agent 研究}

在学术研究层面，LLM4SE 系统综述[1]奠定了该领域的综述性基础，其指出的可控性与评估方法等研究空白至今仍具指导意义。2025 至 2026 年的实证研究将视野从受控实验推进到真实工程环境：文献[2]对十余万 GitHub 项目的大规模分析给出了 coding agent 采用情况的估计；文献[3]基于 AIDev 数据集刻画了 AI 辅助开发活动在真实仓库中的表现，并揭示智能体产出可信度与人类的差距；文献[10]对真实项目中智能体重构（agentic refactoring）的类型、动机与代码质量变化开展了实证研究；相关综述工作亦讨论了智能体在框架、记忆、技能、工具、模型与协作结构等层面的演化，并指出仓库级上下文与可执行反馈是软件工程智能体的重要资源[6]。在评估与系统层面，SWE-bench 构建了基于真实 GitHub issue 的软件工程任务基准[16]，SWE-agent 提出了使语言模型智能体与计算机交互以完成软件工程任务的接口设计[17]，两者共同推动了 coding agent 研究的系统化。总体来看，研究重心正从"模型能否编写代码"转向"如何让智能体在真实工程环境中可靠、可控、低成本地工作"。

在规格驱动开发方向，Piskala（2026）将其概括为从"代码即事实源"（code-as-source）向"规格即首要制品"（specification-as-primary-artifact）的范式转移，并区分了 spec-first、spec-anchored 与 spec-as-source 三种严格等级，且直接分析了 GitHub Spec Kit 等现代工具[7]。

\textbf{（2）项目级规范治理工具：GitHub Spec Kit}

Spec Kit 是 GitHub 开源的规格驱动开发工具包（本文研究时点仓库版本 0.14.x，版本与提交记录固化于附录实验环境），其定位为项目级规范治理与可编排的智能体 SDD 工作流平台[8]。其核心制品包括项目级 Constitution（项目宪法）、功能级 Specification、技术 Plan 与任务分解，并通过 Clarify、Checklist、Analyze、Converge 等质量门禁环节组织需求理解、方案设计、任务拆解、实现与验证，后继版本进一步引入了通用工作流引擎与分层扩展结构。其开发流程同时支持从零构建（greenfield）与存量迭代（brownfield）场景，能力覆盖项目规范、计划生成与任务组织等多个环节，但相应的流程制品与配置体系也带来一定的学习与过程成本。

\textbf{（3）变更级规格管理工具：OpenSpec}

OpenSpec 自 1.0 版本（OPSX Release）起，从固定阶段式工作流转向以动作（Action）为中心的工作流，提出"Actions, not phases"的设计理念，并以制品图（artifact graph，有向无环图）建模制品间依赖与完成状态，强调"依赖是使能条件而非强制关卡"[8]。其以一组围绕变更制品的动作（探索、提案、应用、验证、同步与归档等）组织工作流；在项目组织上强调存量项目优先（specs/ 目录记录系统当前状态，changes/ 目录管理提议变更），并适配了多种主流 AI 编程助手运行环境。其核心定位是变更级规格制品的生命周期管理，而非固定的项目级开发工作流。

\textbf{（4）技能驱动的智能体执行框架：Superpowers}

Superpowers 定位为"面向编程智能体的完整软件开发方法论"（本文研究时点版本 v6.3.0，具体机制随版本迭代，固化于附录实验环境）[9]，以可组合技能与强制工作流组织执行过程：从需求澄清（brainstorming）、Git 工作树隔离、计划编写，到子代理驱动开发（subagent-driven development）、测试驱动开发与代码审查，官方强调"工作流是强制要求而非建议"。其通过技能化工作流、子代理协作与多层评审机制（任务评审、代码质量评审与整分支评审）约束智能体的执行行为，并包含工作区隔离、任务分级等面向过程成本与上下文管理的设计。其治理粒度位于技能与执行工作流层，而非项目级规范宪法体系。

\textbf{（5）国内研究与实践现状}

国内 AI 辅助编程工具以 CodeGeeX、通义灵码、豆包编程助手等为代表，主要聚焦代码补全与单轮问答场景；覆盖需求、设计、计划、测试、编码全流程的工程化工具实践正在起步，与国际开源社区在规格驱动开发与智能体工作流方向上的快速演进相比仍有差距。

\textbf{（6）三者关系与研究空白}

综合来看，Spec Kit、OpenSpec 与 Superpowers 均已具备规格、计划、任务、执行、验证与智能体工作流机制，能力边界重叠度较高（截至本文检索时点，尚未见以三者为对象的同行评审系统比较研究，本文对三者的事实描述以其官方仓库与文档为准），其主要差异体现在设计重心与治理粒度。三者的定位、所解决问题与主要不足对比如表 1-1 所示。

表 1-1 三类代表性工具的定位与不足对比

\begin{table}[htbp]
\centering
\small
\begin{tabularx}{\textwidth}{|X|X|X|X|X|}
\hline
\textbf{工具} & \textbf{治理粒度} & \textbf{核心定位} & \textbf{解决的主要问题} & \textbf{主要不足} \\
\hline
Spec Kit & 项目级 & 规范治理与 SDD 工作流编排 & 规格、计划与任务的流程治理 & 项目知识沉淀与代码知识供给不足 \\
OpenSpec & 变更级 & 规格制品生命周期管理 & 变更提案到归档的规格一致性 & 与代码结构知识的关联较弱 \\
Superpowers & 技能/执行级 & 智能体执行纪律 & 智能体执行行为的约束与评审 & 项目级业务知识的组织较弱 \\
\hline
\end{tabularx}
\end{table}

当前研究与实践分别关注智能体行为控制、规格制品管理或代码知识检索，且各自取得了进展，但尚缺少一种同时考虑项目级知识沉淀、开发流程约束与任务上下文预算优化的统一工程化方法：现有工具与方法通常以文档、制品或智能体工作流为主要组织单元，在"意图-规格-计划-执行"层面已较为成熟，但项目级长期知识（业务能力规范的持续沉淀）与代码结构知识如何进入智能体的任务上下文，以及如何在 Token 预算约束下进行结构感知、可解释的上下文选择，仍缺少面向具体软件项目的工程化方案。这一研究空白构成了本文工作的切入点。


\subsection{代码知识图谱与上下文工程研究现状}
\textbf{（1）研究演进脉络}

代码智能与上下文工程相关研究呈现出清晰的演进脉络。在代码表示学习方向，GraphCodeBERT 显式利用代码的数据流结构辅助代码表示预训练，表明代码表示不应仅依赖 token 序列，代码固有的结构信息能够显著提升下游任务效果[11]；在代码结构建模方向，代码属性图（Code Property Graph）将抽象语法树、控制流图与程序依赖图统一为可查询的图模型，为代码结构知识的图式组织提供了经典范式[18]；在语义代码检索方向，CodeSearchNet 将自然语言到代码的检索确立为研究问题，并发布了约 600 万函数规模的数据集，推动了语义代码检索研究的系统化发展[12]；在需求-代码追踪（Traceability Link Recovery）方向，相关研究长期关注需求、设计与代码之间关联关系的自动恢复问题：Cleland-Huang 等较早提出利用异构信息提升需求追踪收益的方法[13]，Borg 等通过工业案例研究分析了需求到代码链接恢复中的有效性与挑战[14]，后续研究进一步探索了基于机器学习和优化算法的自动化追踪方法[15]；近年来，随着大语言模型的发展，研究者开始探索利用 LLM 提升需求理解、代码语义匹配以及软件知识关联能力，但面向 AI 编程智能体上下文构建场景的需求-代码知识组织仍缺少系统研究。在检索增强与上下文工程方向，RAG 通过引入外部非参数化知识缓解参数化知识的局限[5]，"Lost in the Middle"现象的发现推动了文档压缩与上下文选择研究[4]，GraphRAG 等方法进一步通过在语料上构建实体-关系图并以图结构组织检索与生成过程，探索了利用图结构增强知识检索能力的途径[19]，为本文基于代码知识图谱的上下文生成方法提供了相关启发；面向智能体的上下文组织与选择研究（如智能体记忆机制、仓库级上下文供给等）亦逐渐展开[6]。

\textbf{（2）与 Coding Agent 研究的关联}

值得注意的是，智能体实证研究同样提示了上下文供给的瓶颈地位：文献[3]与文献[10]均显示，智能体参与真实开发时，任务相关上下文的获取方式高度依赖人工检索或整库注入，前者效率低下且易遗漏关键信息，后者则面临 Token 成本与信息稀释的双重压力。

\textbf{（3）国内研究现状}

国内在代码知识图谱与代码表示学习（如基于抽象语法树的代码图谱构建、代码搜索与克隆检测）、需求工程与人工智能结合等方向均有研究开展，但同样以特定任务的学术评测为主。

\textbf{（4）现有局限分析}

综观现有研究，其局限主要体现在三个方面：其一，现有工作较多聚焦于代码表示、代码检索、克隆检测、需求追踪等特定任务，在面向 AI 编程智能体的工程化工作流中，将代码结构知识、业务语义与上下文压缩统一起来的研究仍相对有限；其二，代码知识图谱多以孤立的检索、查询工具形态存在，而非深度嵌入开发工作流的基础设施；其三，针对代码图谱的结构感知压缩与 Token 预算控制研究较少，上下文压缩方法多面向通用文本而非结构化的代码知识。


\section{本文主要工作与创新点}
针对上述研究背景与现状分析，本文围绕"如何让 AI 编程智能体在软件工程中获得项目级知识、受控执行与高效上下文"这一核心问题展开研究。该问题在 1.1.1 节被分解为三个研究问题（RQ1-RQ3），并可进一步归结为一个统一的中间层问题--智能体上下文工程（Agent Context Engineering）：如何为 AI 编程智能体组织、选择、压缩与供给与其当前任务相关的结构化上下文。本文的主要工作即围绕这一问题与三个研究问题展开，凝练为以下三个创新点：

\textbf{创新点 1：面向 AI 编程智能体的流程约束架构}（对应第 4 章）。设计并实现一种面向 AI 编程智能体的三层分离工作流架构（AI 层、CLI 层与文件系统层），以及"CLI 准备-AI 生成交互-CLI 收尾"的三段式编排契约，通过显式的阶段制品（需求文档、能力规格、技术方案、开发计划、测试方案）与 CLI 状态管理机制（状态变更显式留痕、拍板门禁、断点续作）实现 AI 辅助开发的过程约束，使流程确定性控制与模型生成能力协同工作。该创新点解决过程可控性问题（RQ1），对应上下文工程中"上下文在受控流程中被请求与消费"的环节。

\textbf{创新点 2：面向智能体上下文生成的能力-代码双层知识组织方法}（对应第 3 章）。提出 C 层（业务能力）与 L1/L2/L3 层（模块/文件/代码元素）构成的双层图谱模型，将文档提取、命名匹配、语义匹配与 Git 历史等多类启发式证据（以 noisy-OR 概率模型融合）用于自动建立业务-代码关联，降低人工维护需求-代码追踪关系的成本；能力规范随需求归档持续沉淀，形成可供智能体检索的项目级长期知识。需要说明的是，该创新点定位为面向智能体上下文生成场景的知识组织与工程融合方法，而非需求追踪算法层面的创新。该创新点解决业务知识与代码结构脱节问题（RQ2），对应上下文工程中"结构化知识供给"的环节。

\textbf{创新点 3：图谱驱动的任务上下文压缩方法}（对应第 3.3 节）。设计一种图谱驱动的任务上下文压缩方法，涵盖置信度衰减加权的锚点选择、加权双向 BFS 子图裁剪、距离感知的骨架抽取与三档压缩序列化等环节，在 Token 预算约束下按需生成任务相关的高密度代码上下文，减少人工上下文整理成本，并降低大规模仓库整库注入带来的上下文压力。该创新点解决上下文效率问题（RQ3），对应上下文工程中"预算控制"的环节；其有效性以 Token 缩减率、上下文召回率与下游任务表现等指标度量（见第 5 章）。

三个创新点构成"流程约束-知识供给-预算控制"的完整链路。本文设计并实现了 wpw（Web Project Workflow）系统作为上述方法的统一实现与验证平台：第 3 章的知识组织与上下文生成方法、第 4 章的流程约束架构均在该系统中落地，第 5 章在其上开展对比实验与消融评估。其余技术组件（如 noisy-OR 聚合公式细节、BFS 遍历实现、JSONL 与向量存储方案、多语言源码解析器、Plan-as-tracker 任务追踪等）作为支撑上述创新点的实现技术，在相应章节中作为实现细节描述，不单独作为创新点宣称。


\section{文章结构}
本文共分为六章，各章内容安排如下：

第 1 章为绪论。阐述研究背景与意义，综述 AI 辅助开发与智能体工作流编排、代码知识图谱与上下文工程两个方向的国内外研究现状，凝练本文主要工作与三个创新点，并给出文章结构。

第 2 章为相关理论与关键技术介绍。介绍知识图谱的基础概念与本体建模方法、知识表示与嵌入向量模型、大语言模型基础与 Transformer 架构原理，以及检索增强生成与图结构增强检索（RAG/GraphRAG）、智能体架构与上下文工程、需求-代码追踪（Software Traceability）等相关理论，为后续章节的方法设计奠定理论基础。

第 3 章为 wpw 知识库与知识图谱的构建。阐述知识库的数据采集与构建方法、知识图谱的数据处理、多源证据标注、图谱构建与增量更新机制，以及知识库和知识图谱与大语言模型结合的上下文生成方法，涵盖本文创新点 2 与创新点 3。

第 4 章为 wpw 系统设计与实现。从需求分析出发，阐述三层分离架构与 Agent 可靠性设计原则、模块设计、CLI 层与 AI 层实现，以及系统测试，涵盖本文创新点 1。

第 5 章为基于美食推荐系统的对比实验与评估。围绕"不同上下文构建与流程约束策略对 AI 辅助开发效果的影响"设计四方案对比实验与消融实验（wpw 作为其中一种策略实现参与比较），介绍测试项目（基于大模型的美食推荐系统）设计，报告并分析成本、准确性、AI 辅助程度、流程规范性等维度的实验结果，讨论技术栈适用性、优势场景与局限性。

第 6 章为总结与展望。总结本文三个创新点的工作与实验结论，并展望后续研究方向。

\textbf{本章参考文献映射}（正式写作时并入全文参考文献列表，编号待统一）：

\begin{itemize}
\item {}[1] Hou X, et al. Large Language Models for Software Engineering: A Systematic Literature Review. ACM TOSEM, 2024, 33(8).
\item {}[2] Robbes R, et al. Agentic Much? Adoption of Coding Agents on GitHub. arXiv:2601.18341, 2026.
\item {}[3] Li X, et al. The Rise of AI Teammates in Software Engineering (SE 3.0). arXiv:2507.15003, 2025.
\item {}[4] Liu N F, et al. Lost in the Middle: How Language Models Use Long Contexts. arXiv:2307.03172, 2023.
\item {}[5] Lewis P, et al. Retrieval-Augmented Generation for Knowledge-Intensive NLP Tasks. NeurIPS 2020.
\item {}[6] Zhou X, et al. Self-Evolving Coding Agents: A Survey. arXiv:2608.03392, 2026.
\item {}[7] Piskala K. Spec-Driven Development: From Code to Contract in the Age of AI Coding Assistants. arXiv:2602.00180, 2026.
\item {}[8] GitHub. GitHub Spec Kit[EB/OL]. github/spec-kit, 2026.
\item {}[9] OBRA. Superpowers: An agentic skills framework \& software development methodology[EB/OL]. obra/superpowers, v6.3.0, 2026.
\item {}[10] Horikawa Y, et al. Agentic Refactoring: An Empirical Study of AI Coding Agents. arXiv:2511.04824, 2025.
\item {}[11] Guo D, et al. GraphCodeBERT: Pre-training Code Representations with Data Flow. ICLR 2021.
\item {}[12] Husain H, et al. CodeSearchNet Challenge. arXiv:1909.09436, 2019.
\item {}[13] Cleland-Huang J, Zemont G, Lukasik W. A Heterogeneous Solution for Improving the Return on Investment of Requirements Traceability. IEEE International Conference on Requirements Engineering (RE), 2004.
\item {}[14] Borg M, Runeson P, Ardö A. Recovering from Requirements to Code: An Empirical Study of Traceability Link Recovery. Empirical Software Engineering, 2014, 19: 1236-1271.
\item {}[15] Sultanov H, Hayes J H. Application of Genetic Algorithms to the Requirements Traceability Problem. Requirements Engineering, 2010, 15: 173-187.
\item {}[16] Jimenez C E, et al. SWE-bench: Can Language Models Resolve Real-World GitHub Issues? ICLR 2024.
\item {}[17] Yang J, et al. SWE-agent: Agent-Computer Interfaces Enable Automated Software Engineering. NeurIPS 2024.
\item {}[18] Yamaguchi F, Golde N, Arp D, Rieck K. Modeling and Discovering Vulnerabilities with Code Property Graphs. IEEE Symposium on Security and Privacy, 2014.
\item {}[19] Edge D, et al. From Local to Global: A Graph RAG Approach to Query-Focused Summarization. arXiv:2404.16130, 2024.
\end{itemize}
\begin{quote}
全文参考文献规模规划：硕士论文建议 50-80 篇，各方向分布目标--LLM4SE 10 篇+、Coding Agent 8 篇+、RAG/Context Engineering 8 篇+、Software Traceability 5 篇+、知识图谱/代码图谱 10 篇+；后续章节将补充 GraphRAG、Code Property Graph、SWE-bench、Agent Memory 等方向文献，编号在全文合稿时统一。
\end{quote}

